\section{等价关系的原像,诱导的等价关系}

\begin{definition}{等价关系的原像}{preimage-of-equivalent-relation}
	设$E,F$是集合,$f\in\Isom_{\mathsf{Set}}\left(E,F\right)$,$S$是$F$上的等价关系,$u:F\to F/S$是典范投影.映射$u\circ f$在$E$上诱导的等价关系$R$称为$S$在$f$下的原像.
\end{definition}

\begin{definition}{诱导的等价关系}{}
	设$R$是集合$E$上的等价关系,$A\subseteq E$.把$R$在典范单射$A\hookrightarrow E$下的原像记作$R_{A}$,称为$R$在$A$上诱导的等价关系.
\end{definition}

\begin{proposition}{}{}
	设$R$是集合$E$上的等价关系,$A\subseteq E$,则对任意$x\in E$有$\left[x\right]_{R_{A}}=\left[x\right]_{R}\cap A$.
\end{proposition}

\begin{proposition}{}{}
	设$R$是集合$E$上的等价关系,$A\subseteq E$,$p:E\to E/R$和$q:A\to A/R_{A}$是典范投影,典范单射$\iota:A\hookrightarrow E$和$R_{A},R$相容.
	\begin{enumerate}
		\item $\iota$诱导的映射$h:A/R_{A}\to E/R$是单射.
		\item $\im\left(h\right)=f\left(A\right)$.
	\end{enumerate}
	我们称双射$k:A/R_{A}\to f\left(A\right)$及其逆映射为典范双射.
\end{proposition}








